\documentclass[10pt,a4paper]{amsart}
\usepackage[margin=19mm]{geometry}
\usepackage{amsmath,amssymb,mathtools}
\usepackage[scr=rsfs,cal=euler]{mathalfa}
\usepackage{xfrac}
\usepackage[unicode,hidelinks,bookmarks=false]{hyperref}
\newcommand{\ie}{i.e.\ }
\newcommand{\cf}{cf.\ }
\newcommand{\resp}{respectively\ }
\newcommand{\Subsection}{\subsection}
\providecommand{\nobreakdash}{\nobreak\hbox{-}}
\setlength{\emergencystretch}{2em}
\multlinegap0pt
\usepackage{amssymb}
\usepackage{mathtools}
\usepackage{tikz}
\usepackage[shortlabels]{enumitem}
\usepackage{float}
\newtheorem{lem}{Lemma}
\newcommand{\Fp}{\mathscr{F}_\varphi}
\newcommand{\U}{\mathscr{U}}
\title{An Infinite Transitivity Theorem}
\author{Miles Gould}
\date{Source: arXiv:2512.07549v4 (CC BY 4.0)}
\begin{document}
\maketitle

\noindent\emph{Source and licence.} An Infinite Transitivity Theorem. Version arXiv:2512.07549v4 (CC BY 4.0), distributed by arXiv under the Creative Commons Attribution 4.0 International licence, http://creativecommons.org/licenses/by/4.0/, as recorded in the arXiv licence metadata for this exact version and shown on the abstract page https://arxiv.org/abs/2512.07549v4. Full licence notice: \texttt{licenses/arxiv-2512.07549-CC-BY-4.0.txt}.\par\smallskip

\noindent\textit{Editorial scope.} Counted unit: the article's lem fields (source0.tex lines 329--337). The article's complete original proof of it is delivered with it (source0.tex lines 339--362, one \texttt{proof} environment of 24 lines). No mathematics is written, repaired or re-derived: the statement and the proof are the article's own lines, in the article's own notation, and no retained line is substituted or re-flowed.  No further source-local context is needed: every cross-reference of every retained line resolves inside the two delivered spans.  Nothing was omitted from any retained span, so the package carries no omission record.  No retained line contains a citation, so the wrapper carries no bibliography and no bibliography entry.  No retained line contains a figure environment, an image-inclusion command or drawing code, so no image is delivered and the package declares no asset dependency.  Editorial formatting only, in addition: an AMS \texttt{amsart} preamble that restates the article's own declarations and macros used by the retained lines, namely the environments \texttt{lem} on separate counters, together with the standard packages those lines need.  The retained environments print in this document as Lemma 1 in order of appearance, whereas the article prints the same items with the numbering of its own full document.  The complete native source of the article as distributed in the arXiv e-print for this exact version is bundled unchanged as \texttt{sources/190579/source0.tex}.
\begin{lem}\label{fields}
    Let $C$ be countably degree-1 saturated. Then $C$ is saturated over countable types consisting of the following conditions
    \begin{enumerate}
        \item $\|P(\bar{x})\| \in K$
        \item $\|P(\bar{x})\|^\varphi_2 \leq r$
        \item $\varphi(P(\bar{x})) \in \Lambda$
    \end{enumerate}
    where $P(\bar{x})$ a degree-$1$ $*$-polynomial, $K\subseteq[0,\infty)$, $\Lambda\subseteq\mathbb{C}$ both compact, $r\geq 0$, and $\varphi$ a P-state.
\end{lem}

\begin{proof}
    Let $t(\bar{x})=\{\rho_j(P_j(\bar{x}))\in K_j:j\in\mathbb{N}\}$ be a satisfiable countable type consisting of conditions above. To prove that it is realized, we will perform two conversions of $t(\bar{x})$.
    
    Since $t(\bar{x})$ is satisfiable, fix $\bar{b}_m$ such that $\max_{j\leq m}d(\rho_j(P_j(\bar{b}_m)),K_j)\leq\frac{1}{m}$. Then fix an ultrafilter $\U$ on $\mathbb{N}$. Let $s_j=\lim_{m\rightarrow\U}\rho_j(P_j(\bar{b}_m))$. Then
    \[t'(\bar{x})=\{\rho_j(P_j(\bar{x}))=s_j:j\in\mathbb{N}\}\]
    is satisfiable type whose realizations realize $t(\bar{x})$. Now for each $\varphi$, choose $a_\varphi\in\Fp$ excising the $C^*$-subalgebra $B$ generated by coefficients of $P_j(\bar{x})$ and entries of $\bar{b}_m$. Then, for each $j$ such that $\rho_j=\varphi$, choose $d_j\in B$ such that $\varphi(d_j)=s_j$. Define
    \[P'_j(\bar{x}) = \begin{cases} 
        a_\varphi(P_j(\bar{x})-d_j)a_\varphi, & \rho_j=\varphi; \\
        P_j(\bar{x})a_\varphi, & \rho_j=\|\cdot\|_2^\varphi; \\
        P_j(\bar{x}), & \rho_j=\|\cdot\|;
   \end{cases}\]

   \[r_j = \begin{cases} 
        0, & \rho_j=\varphi; \\
        s_j, & \text{otherwise}.
   \end{cases}\]
   
   The type
   \[t''(\bar{x})=\{\|P'_j(\bar{x})\|=r_j:j\in\mathbb{N}\}\]
   is a satisfiable countable degree-1 norm type whose realizations realize $t(\bar{x})$. To see this, we only need to check the case where $\rho_j=\|\cdot\|_2^\varphi$,
   \[\|P_j(\bar{b})\|_2^\varphi=\varphi(P_j(\bar{b})^*P_j(\bar{b}))^\frac{1}{2}\leq\|a_\varphi^*P_j(\bar{b})^*P_j(\bar{b})a_\varphi\|^\frac{1}{2}=\|P_j(\bar{b})a_\varphi\|=r_j\in K_j.\]
   In this case, $K_j$ is of the form $[0,r]$, so $\|P_j(\bar{b})\|^\varphi_2\in K_j$.
   Therefore, since $t''(\bar{x})$ is realized, so is $t(\bar{x})$.
\end{proof}

\end{document}
